\section{Architectural Split Choice}
\label{sec:analysis-split}

\subsection{RQ1.1: Coarse Boundary Selection}
\label{sec:analysis-rq11}

RQ1.1 asks which coarse ResNet-18 stage boundaries provide suitable
two-part microservice decompositions under controlled local execution.
The analytical contribution of this section is not the per-condition
latency itself, which is reported in \cref{tab:rq11-results}, but the
explanation of why the four coarse boundaries are not interchangeable
and which of them is the strongest carry-forward candidate. The
RQ1.1 results support a single coherent reading: the cost of one
service boundary is a joint function of the activation transferred at
that boundary and of how the remaining computation is distributed
between the two services. This matches the framing used in the prior
split-inference literature, which treats partition placement as a
computation--communication trade-off rather than a depth-only
decision~\cite{neurosurgeon,dnnsplit}. The four stage boundaries of
ResNet-18~\cite{resnet} provide a structurally clean coarse sweep, but
the RQ1.1 data make clear that this structural cleanliness is a
necessary, not sufficient, condition for a suitable decomposition.

\paragraph{All four splits are slower than monolithic.}
Under the controlled local conditions used here, no two-part
decomposition improves raw end-to-end latency relative to monolithic
inference: the four splits add between \(+2.4\%\) and \(+4.7\%\) mean
overhead in \cref{tab:rq11-results}. The RQ1.1 result is therefore
diagnostic rather than optimisation-oriented; it identifies the
least-bad boundary and explains the gradient of costs across the
sweep, but it does not support a claim that local two-service
decomposition is a raw-latency improvement. This is consistent with
the broader split-inference literature, where useful latency wins
typically depend on heterogeneous compute or bandwidth conditions that
are deliberately absent from RQ1.1~\cite{neurosurgeon,dnnsplit}.

\paragraph{The earliest boundary is penalised by activation expansion.}
The split after \texttt{layer1} transferred the largest activation
tensor (\(802{,}816\)~B, \(784\)~KiB) and had the largest mean
boundary-crossing interval at \(53.81\)~ms (\cref{tab:rq11-results}).
A pointed comparison is informative: the raw input tensor is
\(1\times3\times224\times224\) FP32, i.e.\ \(602{,}112\)~B, so the
intermediate carried at this boundary is approximately \(1.33\times\)
the raw input. The earliest coarse stage therefore \emph{expands} the
byte budget that the service interface must carry, which is the
cleanest in-thesis explanation for why the earliest split is the most
expensive. Prior work on split inference for CNN backbones reports the
same qualitative pattern --- without representation compression, the
early-stage feature maps can exceed the raw input
size~\cite{bottlenet_pp,deeperthings}. The boundary-crossing interval
reported here is, however, a composite measurement that includes
remote service-B compute and serialisation/deserialisation in addition
to transport (\cref{sec:rq11-results}), so it should not be read as a
pure network-transfer cost.

\paragraph{The latest boundary is penalised by compute degeneracy.}
The split after \texttt{layer4} shows the symmetric failure mode. It
transferred the smallest activation tensor (\(98\)~KiB) and had the
lowest boundary-crossing interval (\(2.22\)~ms), but it left only
\(0.486\)~ms of mean compute on the downstream service, approximately
\(0.62\%\) of total split compute. This is well below the predeclared
\(10\%\) compute-degeneracy threshold (\cref{sec:methods-carryforward})
and the condition is correspondingly flagged in
\texttt{carry\_forward.json}. The implication is methodologically
important: minimising the transferred activation alone can produce a
\emph{structurally weak} microservice decomposition that reduces
transfer burden by collapsing the second service into little more
than the classification head. The same pattern is described in the
distributed-CNN literature, which observes that very late splits
imbalance compute across the partition~\cite{deeperthings,neurosurgeon}.
RQ1.1 therefore retains \texttt{split\_after\_layer4} as evidence for
the transfer-size gradient, not as a credible carry-forward candidate.

\paragraph{The carry-forward region is mid-to-late.}
With the earliest and latest boundaries excluded for the reasons
above, the strongest RQ1.1 candidates are the two mid-to-late
boundaries. The raw-fastest non-degenerate split is
\texttt{split\_after\_layer3} at \(79.08\)~ms, which is selected as
the main carry-forward candidate by the frozen
\texttt{carry\_forward.json}; \texttt{split\_after\_layer2} is
retained as the reference at \(80.12\)~ms. The difference between
these two is modest (approximately \(1.04\)~ms, around \(1.3\%\) of
the slower mean) and should not be read as a general ranking of
\texttt{layer3} over \texttt{layer2} across all environments. The
defensible conclusion is narrower: under RQ1.1's controlled local
conditions, the credible carry-forward region lies around
\texttt{layer2}--\texttt{layer3}, where activation transfer has
already fallen substantially but downstream compute remains
meaningful. This region is what RQ1.2 then refines.

\paragraph{Validation.}
Three pieces of evidence support treating the within-run ordering and
the carry-forward decision as credible for the measured environment.
First, functional parity holds: both local and gRPC parity checks
report \(\texttt{max\_abs\_diff} = 0.0\) at the predeclared tolerance
of \(1.0\times10^{-5}\) across all four split conditions
(\texttt{parity\_validation.json}), so latency differences are not
confounded with numerical drift. Second, round-level behaviour is
stable: the per-condition round coefficient of variation lies between
\(0.002\) (monolithic) and \(0.011\) (\texttt{split\_after\_layer2})
in \texttt{round\_consistency.json}, in line with the
benchmarking-practice guidance to report round-level consistency
rather than mean alone~\cite{benchmarking_sc15}. Third, host controls
(single physical core, thread counts pinned to one, performance
governor, turbo disabled) are applied symmetrically across all
conditions (\cref{tab:rq11-config}), so the comparison is internally
fair. These three legs of the validation argument support the
within-run reading; they do not establish host-independent generality,
because RQ1.1 uses one host, one build, one fixed input seed, CPU-only
batch-one inference, and localhost gRPC, and cross-host or cross-seed
replication is not in scope for this stage.

\paragraph{Evaluation.}
The evaluative reading of RQ1.1 has two facets. Against the
monolithic baseline, every split is measurably slower, with relative
overheads of \(+4.7\%\), \(+3.7\%\), \(+2.4\%\) and \(+3.0\%\) for
splits after \texttt{layer1}, \texttt{layer2}, \texttt{layer3} and
\texttt{layer4} respectively (\cref{tab:rq11-results}). RQ1.1 is
therefore not a result about whether decomposition is faster; it is a
result about which decomposition is least costly while preserving a
meaningful microservice. Within the split set, the carry-forward rule
produces a defensible main--reference pair whose within-run per-iteration
\(95\%\) confidence intervals do not overlap
(\cref{tab:rq11-results}). That separation is a within-run property
and is not a host-independent uncertainty estimate, so it justifies
carrying both \texttt{split\_after\_layer3} and
\texttt{split\_after\_layer2} forward to RQ1.2 rather than committing
to a single boundary at this stage.

\paragraph{Answer to RQ1.1.}
RQ1.1 provides direct but scoped evidence that the most suitable
coarse two-part ResNet-18 microservice decompositions under controlled
local execution are not the earliest or latest stage boundaries. The
earliest boundary is penalised because the layer-1 activation expands
the byte budget the boundary must carry, exceeding the raw input size
in this run. The latest boundary is penalised because it satisfies
the activation-size criterion only by reducing the downstream service
to roughly the classification head, falling below the predeclared
\(10\%\) compute-degeneracy threshold. The strongest carry-forward
candidate from this frozen run is \texttt{split\_after\_layer3}, with
\texttt{split\_after\_layer2} retained as the reference candidate.
Suitability under RQ1.1 conditions is therefore determined by the
\emph{joint} balance of latency, boundary-crossing cost, activation
size and compute distribution, not by any single dimension --- the
two-dimensional reading developed further in
\cref{sec:analysis-rq13}.

\subsection{RQ1.2: Fine-Grained Refinement}
\label{sec:analysis-rq12}

The results in \cref{sec:rq12} support three observations.

\paragraph{The refined region remains tightly clustered.}
All three refined split conditions fall within a very narrow latency band,
from 80.16~ms to 80.39~ms (\cref{tab:rq12-results}). This indicates that the
block-level refinement does not reveal a dramatically superior new split.
Instead, it confirms that the carried-forward \texttt{layer2}--\texttt{layer3}
region identified in RQ1.1 already captures the main structure of the
latency--communication trade-off in this late-stage part of the network.

\paragraph{Equal activation size does not guarantee equal performance.}
The most instructive comparison in this refinement stage is between
\texttt{split\_after\_layer3\_block0} and \texttt{split\_after\_layer3}. Both
conditions transfer an intermediate tensor of approximately 196~KB, yet they
differ in boundary-crossing interval by approximately 7.8~ms and in mean
end-to-end latency by a smaller but still measurable amount. This demonstrates
that activation-transfer size is not the sole determinant of boundary cost. The
distribution of computation on either side of the boundary also matters:
shifting the split from the end of \texttt{layer3} to the internal boundary
after \texttt{layer3.0} changes the compute placement and therefore changes
boundary-crossing behavior even when the serialized tensor is the same size.
This finding is consistent with the broader partitioning literature, which
treats compute distribution as a first-class consideration alongside
communication cost~\cite{dnnsplit, deeperthings, survey_partitioning,
fine_grained_partitioning, hierarchical_partitioning}.

\paragraph{This run does not overturn the coarse screening.}
Under the current refined benchmark outcome, \texttt{split\_after\_layer3\_block0}
is the raw-fastest split and is selected as the main candidate by the
implemented carry-forward rule in
\texttt{rq1\_2\_20260515\_112636/carry\_forward.json}, with
\texttt{split\_after\_layer2} retained as the reference candidate. However,
the three refined conditions remain so close that the central thesis-facing
conclusion is not that RQ1.2 discovered a dramatically better split than
RQ1.1. Rather, the refinement confirms that the late carried-forward region
remains the correct focus, while revealing that small block-level shifts can
still change behavior measurably even when activation-transfer size is held
constant.

\paragraph{Answer to RQ1.2.}
RQ1.2 shows that finer-grained refinement within the carried-forward coarse
region from RQ1.1 does not overturn the coarse-level findings, but it does
reveal a meaningful secondary insight. In the refined comparison,
\texttt{split\_after\_layer3\_block0} is the raw-fastest split and is selected
as the main candidate under the implemented selection rule, while
\texttt{split\_after\_layer2} is retained as the reference candidate. The
refined split \texttt{split\_after\_layer3} is slightly slower in this run.
At the same time, the comparison between \texttt{split\_after\_layer3\_block0}
and \texttt{split\_after\_layer3} --- which share the same activation-transfer
size of approximately 196~KB yet differ in boundary-crossing interval by
approximately 7.8~ms --- demonstrates that equal activation-transfer size does
not guarantee equal performance. Compute placement within the refined region
still matters, even when the serialized tensor is held constant. The coarse RQ1.1
screening had already identified the correct late-stage comparison space; the
block-level refinement confirms that structure and adds the observation that
per-block compute distribution is a secondary but observable factor in
boundary-crossing cost~\cite{dnnsplit, deeperthings, survey_partitioning,
fine_grained_partitioning, hierarchical_partitioning, neurosurgeon, bottlefit}.

\subsection{RQ1.3: Explanatory Synthesis}
\label{sec:analysis-rq13}

The explanatory analysis in \cref{sec:rq13} supports a coherent two-dimensional
account of the local findings from RQ1.1 and RQ1.2.

At the coarse level (RQ1.1), activation-transfer size is the dominant factor
explaining the gradient of boundary-crossing costs from early to late splits
(53.81~ms, 40.85~ms, 26.06~ms, and 2.22~ms after \texttt{layer1} through
\texttt{layer4} respectively, per
\texttt{rq1\_1\_20260515\_111428/condition\_summaries.csv}). The monotonic
relationship between activation size and boundary-crossing interval is strong
and consistent. However, this single dimension is insufficient: the
compute-degenerate \texttt{split\_after\_layer4} minimises transfer burden but
fails as a meaningful decomposition because it collapses one side of the
partition into a near-empty service.

At the refined level (RQ1.2), compute distribution emerges as an independent
explanatory dimension. The controlled comparison between two conditions with
identical activation size but different boundary-crossing intervals shows that
the placement of computation around the boundary affects cost even when
communication burden is held constant. The internal timing experiment supports
this observation by confirming that per-stage and per-block compute within
ResNet-18 is structurally uneven, with \texttt{layer4} accounting for nearly a
third of total forward-pass time and individual blocks within \texttt{layer3}
contributing unequally. The operation-level support data strengthens this
explanation further by showing that the heaviest internal units are
overwhelmingly convolutional and concentrated in the later part of the network.

Neither dimension alone is sufficient. Activation-transfer size explains the
broad coarse pattern but not the refined within-region differences. Compute
distribution explains the refined differences but does not account for the
large boundary-crossing cost gradient between early and late coarse splits.
Together, they provide a more complete explanation: suitable partition
boundaries arise from a trade-off between activation-transfer burden,
boundary-crossing cost, and compute distribution, not from the minimisation of
any single metric.

\paragraph{Answer to RQ1.3.}
RQ1.3 shows that the latency and boundary-crossing behavior observed in RQ1.1
and RQ1.2 can be explained by two complementary structural properties of the
partition: activation-transfer size and compute distribution.

Activation-transfer size is the dominant factor at the coarse level. The
monotonic decrease in boundary-crossing interval from early to late splits in
RQ1.1 closely tracks the reduction in intermediate tensor size as the boundary
moves deeper into the network. This explains why early splits are costly and
why later splits generally incur lower communication overhead.

Compute distribution is an independent and observable factor at the refined
level. The RQ1.2 comparison between \texttt{split\_after\_layer3\_block0} and
\texttt{split\_after\_layer3} --- which share the same activation size of
approximately 196~KB yet differ in boundary-crossing interval by
approximately 7.8~ms --- demonstrates that equal transfer burden does not
guarantee equal performance. The placement of computation around the boundary matters
independently of the size of the transferred tensor.

The internal timing experiment supports the compute-distribution explanation
by confirming that execution cost within ResNet-18 is structurally uneven:
convolution operations account for approximately 82\% of total compute,
\texttt{layer4} alone accounts for 31.28\%, and the heaviest internal units
are concentrated in convolution-heavy late-stage regions
(\cref{tab:rq13-internal-summary,tab:rq13-topops}). These structural
properties explain both the compute degeneracy of the \texttt{layer4} split
observed in RQ1.1 and the boundary-crossing asymmetry between the two 196~KB
conditions observed in RQ1.2.

Taken together, under these controlled local conditions, suitable microservice
boundaries for ResNet-18 arise from a joint trade-off among activation-transfer
burden, boundary-crossing cost, and compute distribution across
services~\cite{neurosurgeon, dnnsplit, jointdnn, deeperthings}. No single
metric is sufficient. The mid-to-late coarse region around \texttt{layer2} and
\texttt{layer3} remains the strongest carry-forward zone because it avoids
both the large activation costs of early boundaries and the poor compute
balance of overly late boundaries. These local explanatory findings provide
the interpretive basis against which the later Kubernetes and Azure stages can
evaluate whether infrastructure context changes the relative importance of
these two dimensions.

\subsection{Synthesis: Architectural Split Choice}
\label{sec:analysis-split-synthesis}

The local split-selection stages show that the position of the model boundary
matters before Kubernetes or cloud effects are introduced. RQ1.1 found that
coarse residual-stage boundaries are not interchangeable: early boundaries
carry high activation-transfer cost, while the very late \texttt{layer4}
boundary becomes compute-degenerate. The strongest carry-forward region is
therefore the mid-to-late part of the network around \texttt{layer2} and
\texttt{layer3}, as shown in \cref{tab:rq11-results}. This matches the broader
split-computing literature, which treats partitioning as a trade-off among
activation size, compute allocation, and deployment
conditions~\cite{neurosurgeon,dnnsplit,deeperthings,survey_partitioning}.

RQ1.2 refined this region and showed that a finer boundary can change behavior
even when activation size is held constant. The comparison between
\texttt{split\_after\_layer3\_block0} and \texttt{split\_after\_layer3} is the
important case: both transfer the same approximate activation size, but their
boundary-crossing intervals differ. This supports the RQ1.3 interpretation
that activation-transfer burden and compute placement are complementary
explanations rather than competing ones~\cite{fine_grained_partitioning,
hierarchical_partitioning,pareto_split}.
